\section{Conditioning of the physical-to-smooth inverse}
\label{sec:v76-conditioning}

The inverse has three different quantitative costs: separating the
physical actions, inverting their actual smooth response, and acquiring
enough successful trajectories.  They depend on different geometric
margins.  We collect them here on the fixed classes used in the preceding
results.  None of the following sufficient bounds is a lower bound for
all possible reconstruction procedures.

\subsection{Contraction, differentiability and phase weights}

For the scalar visit bound $a<1$, the sufficient weighted threshold and
its least integer order are exactly
\begin{equation}\label{eq:v76-order}
 \theta_m<1\quad\Longleftrightarrow\quad a^m<\frac17,
 \qquad m_{\min}=\max\left\{3,
       \left\lfloor\frac{\log7}{-\log a}\right\rfloor+1\right\}.
\end{equation}
At a chosen order the weighted inverse prefactor is
$2/(1-\theta_m)$.  Minimizing the order need not give a comfortable
inverse margin.  For any chosen $0<\vartheta<1$, the choice
\begin{equation}\label{eq:v76-order-margin}
 m\ge\max\left\{3,
  \left\lceil\frac{\log((6+\vartheta)/\vartheta)}{-\log a}
                   \right\rceil\right\}
\end{equation}
ensures $\theta_m\le\vartheta$ and hence a prefactor at most
$2/(1-\vartheta)$.  These formulas follow by solving
$6a^m/(1-a^m)\le\vartheta$; the strict inequality in
\eqref{eq:v76-order} explains the floor followed by $+1$.

For example, the scalar bounds give the following values.  The last
column uses $s=1$ in the one-dimensional smoothing exponent; it is not
an asserted optimal statistical rate.
\begin{center}
\begin{tabular}{ccccc}
\toprule
$a$ & $m_{\min}$ & $\theta_{m_{\min}}$ &
 $2/(1-\theta_{m_{\min}})$ & $\alpha_1$\\
\midrule
$1/2$ & $3$ & $0.857143$ & $14$ & $1/9$\\
$9/10$ & $19$ & $0.937099$ & $31.7961$ & $1/41$\\
$99/100$ & $194$ & $0.995514$ & $445.8342$ & $1/391$\\
\bottomrule
\end{tabular}
\end{center}
Taking instead $\vartheta=1/2$ gives orders $4,25,256$, respectively,
and weighted prefactor at most $4$.  It increases the differentiability
required of the observed action and the graph prior.  The jet-alignment
and forward constants also depend on this order; bounding the weighted
prefactor alone does not bound the entire inverse.

Corollaries~\ref{cor:v76-phase-separation} and
\ref{cor:v76-phase-stability} replace $a$ in these sufficient tests by
$\overline a=(\prod_i\widehat a_i)^{1/r}$ when phase-specific bounds
are available on the enlarged class.  The unweighted prefactor then
contains $\kappa(w)$.  The product of the phasewise upper bounds must
not be replaced by the supremum of the products along one orbit without
a new uniform estimate.  Likewise an exactly computed linear multiplier
does not itself provide the nonlinear collar bounds: the latter follow
by continuity on the stipulated class.  The scalar criterion remains
valid and is often simpler to use.

\subsection{Observation margins}

In the two-offset quotient, put $h=A(u)+C(v)$ and
$\Delta=d_2-d_1>0$.  On the exact image,
\[
 d_2-d_1Q=\frac{d_2\Delta}{d_2-h},\qquad
 \left|\frac{\partial H}{\partial Q}\right|
        =\frac{d_1(d_2-h)^2}{d_2\Delta}.
\]
Thus offset separation and the bounds for the positive residual enter
the action inverse explicitly.  Positive density floors are also
required to form $Q$; its higher derivatives involve the same reciprocal
functions.  This route does not require a nonzero obliquity $p$ and
therefore retains the normal-incidence case.

For a single law, the inverse additionally uses
$\kappa_f=p^2/d^2$ and the fixed-slice anchors $u_a,v_a$ of
\eqref{eq:v72-anchors}.  The map $\kappa\mapsto|p|/\sqrt\kappa$
has derivative of absolute value $|p|/(2\kappa^{3/2})$.
The geometric bounds after Theorem~\ref{thm:v72-separation} supply
$\kappa_f\ge p_*^2/d_+^2$ and
$|u_a|,|v_a|\ge p_*a/(2D_*)$.  Their degeneration is independent
of the order condition for the smooth envelope.

Moment reconstruction brings its own matrix condition.  In the affine
model of Remark~\ref{rem:v75-conditioning}, put
$z=pR^2/(3d)$.  Then
\[
 H=\begin{pmatrix}1&-z\\z&0\end{pmatrix},\qquad
 \det H=z^2,\qquad
 H^{-1}=\begin{pmatrix}0&1/z\\-1/z&1/z^2\end{pmatrix}.
\]
The shrinking-gate and small-obliquity losses are visible in the inverse,
not just its determinant.  This is an algebraic density example, not an
asserted positive-curvature billiard realization.  Rescaling the moments
changes their units and the corresponding error norms; it does not
remove the physical conditioning from a comparison in fixed units.

These requirements are separate from the positive Jacobi mass,
nongrazing and clearance bounds that construct the stationary trajectories.
They are also separate from the uniform $C^{m+3}$ graph prior used to
align jets.  On a fixed class with all these margins, their contributions
are finite; uniformity is not asserted as any of them degenerates.

\subsection{Preparation budget and the exact finite anchor}

For the moment protocol the conditional-sample and bounded-readout
exponents are
\[
 \alpha_1=\frac{s}{2(m+s)+1},\qquad
 \gamma_1=\frac{s}{m+s+2}.
\]
They arise from the complete bound in
Lemma~\ref{lem:v74-estimation}, not just its square-root term.
In particular the Bernstein envelope, interpolation-grid and
post-acceptance perturbation terms remain part of the estimate.
Suppose the uniform finite-flight defect is $Ce^{-\omega N}$ and
the success probability is at least $c e^{-\Gamma N}$.
For $B$ independent preparations per phase, the two terms to balance
have the form
\[
 e^{-\omega N},\qquad
 \left(\frac{L_Be^{\Gamma N}}B\right)^{\alpha_1}.
\]
Solving their equality gives
\begin{equation}\label{eq:v76-budget-balance}
 N\sim\frac{\alpha_1}{\omega+\alpha_1\Gamma}
                      \log\frac{B}{L_B},\qquad
 \beta_1=\frac{\alpha_1\omega}{\omega+\alpha_1\Gamma}.
\end{equation}
Rounding to a returning number of flights changes $N$ by a bounded
amount.  The count, bandwidth, minimum-flight and small-error conditions
in Theorem~\ref{thm:v74-charged} and
Corollary~\ref{cor:v74-budget} must still hold.  Increasing $m$ may
improve the weighted inverse margin while decreasing $\alpha_1$;
rarity decreases it further to $\beta_1$.  A positive exponent is not
a bound for the computational cost of finding the selected actual table.

Counts enter area recovery through the exact finite identity
\[
 \mathcal A(T)=\frac{D_{N,b}(c)d_b}
                    {2\pi\pi_{N,b}f_{N,b}(0,0)}.
\]
The count and the fitted actual finite density use the same full gate.
Replacing the latter by its rank-two limit would discard a finite-flight
term.  The logarithmic derivative of $D_{N,b}$ costs $1+N$, as shown
in Section~\ref{sec:v73-complete-record}; at
\eqref{eq:v76-budget-balance} it produces the retained logarithmic area
loss.  Failures are charged, gate and itinerary tags are exact, and
readout perturbations occur after acceptance.  The independently reset
Liouville experiment, not an arbitrary unknown preparation measure,
supplies the normalizer.  These are conditions on the stated inverse
problem, not additional conclusions of the conditioning calculation.
